\chapter{Integridade Acadêmica e Uso de Tecnologias}
\label{chap:integridade}

A integridade acadêmica sustenta a confiança na produção universitária. Ela envolve honestidade intelectual, respeito às fontes, autoria responsável, transparência metodológica e compromisso com a verdade possível dentro dos limites da pesquisa.

\section{Plágio}

Plágio ocorre quando ideias, textos, dados, imagens, estruturas ou resultados de outra pessoa são apresentados como se fossem próprios. Pode ocorrer por cópia literal sem citação, paráfrase sem referência, tradução sem indicação, uso de trabalho alheio ou reaproveitamento indevido de texto próprio sem transparência.

Mesmo quando não há intenção de enganar, descuido com citações e referências pode gerar problema grave. Por isso, o estudante deve registrar fontes desde o início da pesquisa.

\section{Tipos comuns de plágio}

\begin{longtable}{p{0.28\linewidth}p{0.60\linewidth}}
\toprule
Tipo & Caracterização \\
\midrule
\endfirsthead
\toprule
Tipo & Caracterização \\
\midrule
\endhead
Cópia literal & Reprodução de trecho sem aspas, destaque ou indicação correta da fonte. \\
Paráfrase sem fonte & Reescrita de ideia de outro autor sem citação. \\
Plágio mosaico & Combinação de trechos de várias fontes com pequenas alterações, sem autoria clara. \\
Tradução sem indicação & Tradução de texto estrangeiro apresentada como redação própria. \\
Uso indevido de imagem ou tabela & Inserção de figura, tabela, gráfico ou esquema sem fonte ou autorização quando necessária. \\
Autoplágio & Reaproveitamento de texto próprio já avaliado ou publicado sem transparência. \\
Referência falsa & Indicação de fonte inexistente, não consultada ou que não sustenta a afirmação feita. \\
Terceirização de autoria & Entrega de texto produzido por outra pessoa ou ferramenta como se fosse autoria do estudante. \\
\bottomrule
\end{longtable}

\section{Como prevenir}

Algumas práticas reduzem o risco de plágio:

\begin{itemize}
  \item registrar a fonte completa no momento da leitura;
  \item anotar páginas de trechos importantes;
  \item separar claramente citação literal, paráfrase e comentário próprio;
  \item usar aspas ou bloco próprio em transcrições literais;
  \item revisar se toda ideia externa foi citada;
  \item conferir se todas as referências correspondem a fontes realmente consultadas.
\end{itemize}

\section{Autoplágio e reaproveitamento}

O estudante pode retomar trabalhos anteriores quando isso for pertinente e autorizado, mas deve informar a origem e evitar apresentar produção já avaliada como se fosse nova. O reaproveitamento precisa ser transparente.

\section{Fabricação e manipulação de dados}

Inventar participantes, alterar respostas, omitir resultados relevantes, selecionar apenas dados convenientes ou apresentar análise que não foi realizada compromete a validade do trabalho.

Resultados inesperados ou negativos também são resultados. O TCC deve relatar o percurso com honestidade.

Fabricação é a criação de dados, participantes, fontes ou resultados inexistentes. Falsificação é a alteração, omissão ou manipulação indevida de materiais, procedimentos ou resultados. Ambas comprometem a validade acadêmica, ainda que o texto esteja formalmente bem apresentado.

\section{Uso de inteligência artificial}

Ferramentas de inteligência artificial podem apoiar revisão linguística, geração inicial de ideias, organização, tradução, síntese, programação, criação de imagens e exploração de dados, desde que seu uso respeite as regras do curso e não substitua a autoria do estudante. O uso permitido depende da finalidade, da disciplina e do regulamento aplicável.

\begin{longtable}{p{0.30\linewidth}p{0.58\linewidth}}
\toprule
Uso & Verificação indispensável \\
\midrule
Revisão e tradução & Conferir sentido, terminologia, voz autoral e alterações de conteúdo. \\
Ideias e sínteses & Retomar as fontes originais e reconstruir o argumento com compreensão própria. \\
Programação & Testar o código, licenças, segurança, resultados e adequação metodológica. \\
Imagens & Declarar a geração quando exigido e verificar direitos, vieses e integridade visual. \\
Análise de dados & Validar método, parâmetros, código, resultados e possibilidade de reprodução. \\
Busca bibliográfica & Confirmar cada obra em catálogo, base ou fonte confiável; a ferramenta não substitui a busca. \\
Referências & Nunca aceitar referência sem verificar autoria, título, veículo, data, DOI ou endereço. \\
Produção de texto & Preservar autoria intelectual, transparência e responsabilidade; texto gerado não se torna automaticamente produção do estudante. \\
\bottomrule
\end{longtable}

O estudante não deve entregar texto produzido por ferramenta como se fosse produção própria sem revisão, compreensão e responsabilidade. Também não deve inserir referências inexistentes, dados inventados ou interpretações não verificadas.

Outros riscos incluem falsa citação, reprodução de vieses, exposição de dados pessoais ou sigilosos, violação de direitos autorais, geração de imagens enganosas, perda de confidencialidade e dependência de resultados que o estudante não consegue explicar.

Quando o curso exigir declaração de uso de ferramentas, o estudante deve descrevê-la de forma transparente. O Apêndice \ref{ap:declaracao-ia} apresenta um modelo de apoio do projeto, cuja presença não cria obrigatoriedade. Ferramentas de IA não são autoras do trabalho e não assumem responsabilidade acadêmica.

\fluxopendente

A exigência da declaração, sua posição no trabalho, o fluxo de análise e eventuais consequências dependem da política institucional e das regras do curso.

\section{Similaridade textual e avaliação de integridade}

Relatórios de similaridade são instrumentos de triagem. Um percentual, isoladamente, não prova plágio e tampouco assegura integridade. A análise deve considerar a origem das coincidências, o uso de citações, a natureza do documento, as práticas da área e a presença de apropriação indevida.

\section{Direitos autorais e licenças}

Citar a fonte é indispensável, mas nem sempre basta para autorizar a reprodução integral de imagem, instrumento, escala, mapa, fotografia, software ou outro material protegido. O estudante deve observar a Lei nº 9.610/1998, licenças de uso, termos das plataformas e autorizações necessárias. Materiais em domínio público ou com licença aberta também exigem atribuição conforme as condições aplicáveis.

\section{Responsabilidade sobre a versão final}

Independentemente das ferramentas utilizadas, o estudante é responsável por todo o conteúdo entregue: afirmações, citações, referências, dados, imagens, tabelas, conclusões e normalização.

\caixaorientacao{Uma prática simples reduz riscos: mantenha uma lista de fontes consultadas, registre páginas importantes e nunca copie trechos para o rascunho sem marcar claramente que são citações.}
